\begin{statementbox}
\textbf{C016\quad 双曲线焦点三角形四心坐标性质}

设双曲线
\[
C:\frac{x^2}{a^2}-\frac{y^2}{b^2}=1,\qquad
 a>0,\ b>0,\quad c^2=a^2+b^2,\quad e=\frac ca>1,
\]
左右焦点为 $F_1(-c,0)$、$F_2(c,0)$，$P(x_0,y_0)$ 是双曲线上\textbf{非顶点}的动点。
此时 $|x_0|>a$、$y_0\ne0$，三角形 $PF_1F_2$ 不退化。
记 $\varepsilon=\operatorname{sgn}(x_0)$（右支取 $1$，左支取 $-1$），则：
\begin{enumerate}
  \item \textbf{重心}：
  \[\boxed{G\left(\frac{x_0}{3},\frac{y_0}{3}\right)}.\]
  \item \textbf{内心}：
  \[
  \boxed{I\left(\varepsilon a,\frac{cy_0}{c+e|x_0|}\right)
  =\left(\varepsilon a,\frac{ay_0}{a+|x_0|}\right)}.
  \]
  特别地，右支时 $I\left(a,\frac{ay_0}{a+x_0}\right)$；
  左支时 $I\left(-a,\frac{ay_0}{a-x_0}\right)$。
  \item \textbf{垂心}：
  \[\boxed{H\left(x_0,\frac{c^2-x_0^2}{y_0}\right)}.\]
  \item \textbf{外心}：
  \[\boxed{O'\left(0,\frac{x_0^2+y_0^2-c^2}{2y_0}\right)}.\]
\end{enumerate}

\textbf{特别重要的轨迹结论：}
内心始终位于两条顶点切线 $x=\pm a$ 上，而且
\[
\boxed{x_I\in\{-a,a\},\qquad 0<|y_I|<b.}
\]
因此，\textbf{内心的准确轨迹}是 $x=a$ 和 $x=-a$ 上满足
$0<|y|<b$ 的四条开线段，不能把两条完整直线当作轨迹。

\textbf{通用几何结论：}
对任意不退化的焦点三角形（无需预先知道 $P$ 在什么曲线上），
\[
\boxed{x_I=\frac{|PF_1|-|PF_2|}{2}},\qquad
\boxed{y_I=\frac{c\,y_0}{c+\bigl(|PF_1|+|PF_2|\bigr)/2}}.
\]
双曲线定义保证两焦半径之差恒为 $\pm2a$，故内心横坐标恒为 $\pm a$。
\end{statementbox}
